% =====================================================================
%  Style — applique les directives Sayadi
% =====================================================================

% ----- Interligne 1.5 ------------------------------------------------
\onehalfspacing

% ----- En-tête / pied de page ----------------------------------------
% Hauteur d'en-tête suffisante pour le texte 12 pt (évite les avertissements
% « \headheight is too small » de fancyhdr).
\setlength{\headheight}{15pt}
\pagestyle{fancy}
\fancyhf{}
\fancyhead[L]{\nouppercase{\leftmark}}
\fancyhead[R]{\thepage}
\renewcommand{\headrulewidth}{0.4pt}
\renewcommand{\footrulewidth}{0pt}

% Page de premier chapitre — style « plain » sans en-tête
\fancypagestyle{plain}{%
  \fancyhf{}%
  \fancyfoot[C]{\thepage}%
  \renewcommand{\headrulewidth}{0pt}%
}

% ----- Mise en forme des titres de chapitre --------------------------
\usepackage{titlesec}
% « Chapitre N : Titre » sur une seule ligne (format hang, et non display)
\titleformat{\chapter}[hang]
  {\normalfont\LARGE\bfseries}
  {\chaptertitlename\ \thechapter\,:}{0.5em}{}
\titlespacing*{\chapter}{0pt}{0pt}{24pt}

\titleformat{\section}{\normalfont\large\bfseries}{\thesection}{1em}{}
\titleformat{\subsection}{\normalfont\normalsize\bfseries}{\thesubsection}{1em}{}

% ----- Justification du texte ----------------------------------------
% Justifié par défaut dans `report` ; on désactive l'over/underfull warnings cosmétiques
\tolerance=1000
\emergencystretch=1em

% ----- Profondeur du sommaire ----------------------------------------
\setcounter{tocdepth}{3}
\setcounter{secnumdepth}{3}
